\chapter{Introduction}
\label{ch:introduction}

\section{Motivation and Problem Statement}
\label{sec:motivation}

\todo{Alps Alpine's dependency on non-EU hyperscalers. The Local Products
Department is exploring alternatives to public AI services for data-security
reasons. Frame the concrete business problem before the regulatory one.}

\section{Research Questions}
\label{sec:research-questions}

\todo{State RQ1--RQ3 explicitly. Each must be answerable by the studies in
\cref{ch:evaluation,ch:build} and is answered in \cref{sec:rq-responses}.}

\begin{description}
  \item[RQ1] \todo{Can EU cloud providers satisfy the EC DG DIGIT Cloud
        Sovereignty Framework requirements?}
  \item[RQ2] \todo{How does the readiness of EU providers compare to AWS?}
  \item[RQ3] \todo{What is the trade-off between sovereignty and readiness?}
\end{description}

\section{Scope and Delimitation}
\label{sec:scope}

\todo{State plainly what is NOT covered. Most importantly: no live deployment
is performed — the build branch is assessed analytically as a feasibility
study. Also: single instance tier, point-in-time assessment, requirements
drawn from one enterprise.}

\section{Structure of this Thesis}
\label{sec:structure}

\todo{One short paragraph mapping the six chapters. Note the build-vs-buy
split: \cref{ch:evaluation} answers "which provider do we buy from",
\cref{ch:build} answers "what would it take to build it ourselves".}
